% ============================================================
%  part1/ch06-embeddings.tex  —  Chapter 6: Embeddings
% ============================================================

\chapter{Machine Learning Embeddings}
\label{ch:embeddings}

\epigraph{%
  The purpose of compression is not to reduce\\
  but to reveal.%
}{---}

\begin{WhyMatters}
  Embeddings are the contemporary form of a very old problem:
  how to represent a complex object in a simpler space while
  preserving what matters.
  This chapter bridges the human examples of earlier chapters
  and the cosmological example of Chapter~\ref{ch:horizons}.
  It also introduces, informally, the adjunction $F \dashv G$
  that will be made precise in Chapter~\ref{ch:analysis-synthesis}.
\end{WhyMatters}

\section{Embeddings as controlled compression}

Let $X$ be a space of objects: texts, images, sounds, documents,
observations, or states.
A machine learning embedding is a map
\[
  E : X \to \mathbb{R}^n.
\]
It replaces an object $x \in X$ with a vector $E(x)$.

The embedding is useful when relevant structure in $X$ is
approximately preserved by distances in $\mathbb{R}^n$.
If $d_X$ is a task-relevant distance on $X$, define the
embedded distance
\[
  d_E(x, y) = \|E(x) - E(y)\|.
\]
A good embedding satisfies, approximately,
\[
  d_X(x, y) \approx d_E(x, y)
\]
for the distinctions relevant to the task.

\begin{definition}[Embedding equivalence]
  Two objects $x, y \in X$ are \emph{embedding-equivalent}
  when $E(x) = E(y)$.
  We write $x \sim_E y$.
\end{definition}

\begin{proposition}
  The embedding map $E : X \to \mathbb{R}^n$ induces a quotient
  space $X / {\sim_E}$.
\end{proposition}

\begin{proof}
  Equality of images is an equivalence relation.
  The quotient exists as the set of equivalence classes.
\end{proof}

\begin{Misconception}
  Higher-dimensional embeddings do not necessarily preserve more
  meaning.
  An embedding into $\mathbb{R}^{10000}$ may discard more
  task-relevant structure than an embedding into $\mathbb{R}^{32}$
  if the higher-dimensional space encodes irrelevant variation.
  Admissibility of an embedding depends on the task, not on
  the dimension of the target space.
\end{Misconception}

\section{The inverse image and reconstruction}

Given an embedded vector $v \in \mathbb{R}^n$, the compatible
original objects are
\[
  E^{-1}(v) = \{x \in X : E(x) = v\}.
\]

\begin{proposition}
  If there exists $v \in \mathbb{R}^n$ with $|E^{-1}(v)| > 1$,
  then no reconstruction map $G : \mathbb{R}^n \to X$ can
  recover every $x \in X$ exactly.
\end{proposition}

\begin{proof}
  If $x_1 \neq x_2$ with $E(x_1) = E(x_2) = v$, then any
  $G$ satisfying $G(E(x_i)) = x_i$ must satisfy both
  $G(v) = x_1$ and $G(v) = x_2$, forcing $x_1 = x_2$.
  Contradiction.
\end{proof}

The reconstruction problem is therefore not $E^{-1}$
(which usually does not exist as a function) but
\[
  G : \mathbb{R}^n \to \hat{X},
\]
where $G(v)$ selects an admissible reconstruction from
$E^{-1}(v)$.

\section{Task-relevance and admissible compression}

\begin{definition}[Admissible embedding]
  An embedding $E : X \to \mathbb{R}^n$ is
  \emph{admissible for a task} $R$ when the reconstruction
  error remains below a task threshold:
  \[
    d_R(x, G(E(x))) < \varepsilon_R
  \]
  for all task-relevant $x \in X$.
\end{definition}

The best embedding is not necessarily the one that preserves
the most information:
\[
  \arg\max_E A_R(E)
  \neq
  \arg\max_E |I(E)|.
\]
This is the microscope principle applied to representation
design.

\begin{example}[Word embeddings]
  A word embedding $E : \mathrm{Vocabulary} \to \mathbb{R}^{300}$
  maps words to vectors so that semantically similar words
  map to nearby vectors.
  The embedding discards spelling, etymology, and phonetic
  structure while preserving distributional meaning.
  This is admissible for tasks requiring semantic similarity
  and inadmissible for tasks requiring morphological analysis.
\end{example}

\section{The adjunction, informally}

The pair $(E, G)$ is the first informal appearance of an
adjunction $F \dashv G$.

Encoding $E$ decomposes objects into vectors.
Reconstruction $G$ reassembles vectors into objects.
They are not inverses: $G \circ E \neq \mathrm{id}_X$ in
general.
Instead, there is a systematic relationship: maps from encodings
to targets correspond to maps from reconstructions to originals.

\begin{ForwardRef}
  The informal adjunction $(E, G)$ introduced here will be
  made precise in Chapter~\ref{ch:analysis-synthesis} as
  $F \dashv G : \mathbf{Loc} \to \mathbf{Glob}$.
  The unit $\eta : \mathrm{id} \to G \circ F$ measures
  how faithfully the Monoturn is recovered after decomposition
  and reconstruction.
  The counit $\varepsilon : F \circ G \to \mathrm{id}$
  measures how well a proposed local dataset is explained by
  a reconstructed whole.
\end{ForwardRef}

\section{From embeddings to horizons}

Some projections are chosen. Others are unavoidable.

A cartographer chooses a map projection.
An engineer chooses an embedding space.
An observer does not choose a horizon.
Yet mathematically all three confront the same problem:
reconstruction from incomplete access.

The embedding map $E : X \to \mathbb{R}^n$ is designed.
The horizon restriction
$F_O : \mathcal{M} \to \mathcal{A}(H_O)$
is imposed by physical accessibility.

Both create the same formal question:
\[
  \text{Given non-injective access, what can be reconstructed?}
\]

\begin{MathEcho}
  This chapter introduced:
  \[
    E : X \to \mathbb{R}^n,
    \quad
    x \sim_E y,
    \quad
    E^{-1}(v),
    \quad
    G : \mathbb{R}^n \to \hat{X},
    \quad
    d_R(x, G(E(x))) < \varepsilon_R.
  \]
  The informal adjunction $(E, G)$ anticipates
  $F \dashv G$ in Chapter~\ref{ch:analysis-synthesis}.
  The admissible embedding definition anticipates the
  admissibility sheaf in Chapter~\ref{ch:sheaf}.
\end{MathEcho}

\WhatSurvived{Embedding $E : X \to \mathbb{R}^n$}{%
  Task-relevant similarity\\
  Approximate metric structure\\
  Linear combinations of features%
}{%
  Exact identity of objects\\
  Task-irrelevant variation\\
  Objects outside training distribution%
}{%
  Partial; $G(E(x))$ is admissible\\
  but not exact in general.%
}{%
  $\ker(E)$: all pairs $(x_1,x_2)$\\
  with $E(x_1)=E(x_2)$;\\
  typically dense in $X$.%
}
